% =========================================================
% Open Sets, Closed Sets, Interior, Exterior, Boundary
% Source: volume-iii/analysis/metric-spaces/notes/open-closed-sets.tex
%
% Dependencies:
%   notes-balls.tex          — def:open-ball, def:closed-ball,
%                              rem:interior-ball-identity (radial gap)
%   notes-working-with-balls.tex — prop:reverse-triangle
% =========================================================

% ---------------------------------------------------------
% Toolkit
% ---------------------------------------------------------

\begin{tcolorbox}[colback=gray!6, colframe=gray!40, arc=2pt,
  left=6pt, right=6pt, top=4pt, bottom=4pt,
  title={\small\textbf{Open Sets, Closed Sets, Interior, Exterior, Boundary --- Quick Reference}},
  fonttitle=\small\bfseries]
\small
\begin{tabular}{l l l l}
\toprule
\textbf{Label} & \textbf{Statement} & \textbf{Proof method} & \textbf{Detail} \\
\midrule
D\ref{def:open-closed-set}
  & $S$ open $\iff \forall x\in S,\;\exists r>0,\;B_r(x)\subseteq S$
  & Definition
  & \hyperref[def:open-closed-set]{$\downarrow$} \\[4pt]
P\ref{prop:trivial-open-closed}
  & $\varnothing$ and $X$ are open and closed
  & Vacuous truth; complement
  & \hyperref[prop:trivial-open-closed]{$\downarrow$} \\[4pt]
P\ref{prop:balls-open-closed}
  & $B_r(x)$ is open;\;\, $\overline{B}_r(x)$ is closed
  & Radial gap + triangle inequality
  & \hyperref[prop:balls-open-closed]{$\downarrow$} \\[4pt]
D\ref{def:interior}
  & $\operatorname{int}(A) = \{x\in A : \exists r>0,\;B_r(x)\subseteq A\}$
  & Definition
  & \hyperref[def:interior]{$\downarrow$} \\[4pt]
D\ref{def:exterior}
  & $\operatorname{ext}(A) = \operatorname{int}(X\setminus A)$
  & Definition
  & \hyperref[def:exterior]{$\downarrow$} \\[4pt]
D\ref{def:boundary}
  & $x\in\partial A \iff \forall r>0$: $B_r(x)$ meets both $A$ and $X\setminus A$
  & Definition
  & \hyperref[def:boundary]{$\downarrow$} \\[4pt]
T\ref{thm:ibe-decomp}
  & $X = \operatorname{int}(A)\;\dot\cup\;\partial A\;\dot\cup\;\operatorname{ext}(A)$
  & Set inclusion both ways
  & \hyperref[thm:ibe-decomp]{$\downarrow$} \\[4pt]
T\ref{thm:boundary-closure}
  & $\partial A = \overline{A}\cap\overline{(X\setminus A)}$;\quad $\overline{A}=A\cup\partial A$
  & Closure and boundary definitions
  & \hyperref[thm:boundary-closure]{$\downarrow$} \\[4pt]
P\ref{prop:closed-owns-boundary}
  & $F$ closed $\iff \partial F \subseteq F$
  & Complement + boundary definition
  & \hyperref[prop:closed-owns-boundary]{$\downarrow$} \\[4pt]
\bottomrule
\end{tabular}
\end{tcolorbox}

% ---------------------------------------------------------
\subsection{Definitions: Open Sets and Closed Sets}
\label{sec:open-closed-definitions}
% ---------------------------------------------------------

\begin{tcolorbox}[colback=propbox, colframe=propborder, arc=2pt,
  left=6pt, right=6pt, top=4pt, bottom=4pt,
  title={\small\textbf{Definition (Open Set, Closed Set)}},
  fonttitle=\small\bfseries]
\begin{definition}[Open Set, Closed Set]
\label{def:open-closed-set}
Let $(X,d)$ be a metric space and let $S \subseteq X$.
\begin{itemize}
  \item $S$ is \emph{open} (in $X$) if for every $x \in S$ there exists
        $r > 0$ such that $B_r(x) \subseteq S$.
  \item $S$ is \emph{closed} (in $X$) if its complement $X \setminus S$
        is open.
\end{itemize}
\end{definition}
\end{tcolorbox}

\begin{remark*}[Fully quantified form]
\begin{align*}
S\ \text{open}
  &\iff \forall x \in S,\;
         \exists r \in \mathbb{R}^{+},\;
         B_r(x) \subseteq S. \\[4pt]
S\ \text{closed}
  &\iff X \setminus S\ \text{is open} \\
  &\iff \forall x \in X \setminus S,\;
         \exists r \in \mathbb{R}^{+},\;
         B_r(x) \subseteq X \setminus S.
\end{align*}
\end{remark*}

\begin{remark*}[Intuition]
Every point of an open set has room around it: no point of $S$ is
pressed against the edge of $S$.
Closed sets are defined indirectly as complements of open sets; the
intuition that closed sets contain their limits is a consequence,
not the definition.
See \hyperref[sec:sequential-characterizations]{\S\ref*{sec:sequential-characterizations}}.
\end{remark*}

\begin{remark*}[Common error]
Open and closed are \emph{not} opposites.
A set can be both open and closed (called \emph{clopen}), or neither.
In $\mathbb{R}$: the sets $\mathbb{R}$ and $\varnothing$ are clopen;
the half-open interval $[0,1)$ is neither open nor closed.
\end{remark*}

% =========================================================
% TikZ diagram: Open set definition
% Place after the Open Set / Closed Set definition box.
% Requires: \usepackage{tikz}  \usetikzlibrary{calc}
%           \usepackage{caption}   (for \captionof)
% =========================================================

\begin{center}
\begin{tikzpicture}[
  dot/.style     = {circle, fill, inner sep=1.8pt},
  lbl/.style     = {font=\small},
  sublbl/.style  = {font=\scriptsize, text=gray!55!black}
]

% --- Open set S ---
\fill[blue!10]
  (0, -0.3)
  .. controls ( 1.4,  0.2) and ( 2.2,  1.4) .. ( 1.8,  2.8)
  .. controls ( 1.2,  4.0) and (-0.6,  4.4) .. (-1.8,  3.6)
  .. controls (-3.2,  2.8) and (-3.4,  1.0) .. (-2.4,  0.0)
  .. controls (-1.6, -0.8) and (-0.6, -0.6) .. ( 0,   -0.3)
  -- cycle;

\draw[blue!55!black, dashed, line width=0.9pt]
  (0, -0.3)
  .. controls ( 1.4,  0.2) and ( 2.2,  1.4) .. ( 1.8,  2.8)
  .. controls ( 1.2,  4.0) and (-0.6,  4.4) .. (-1.8,  3.6)
  .. controls (-3.2,  2.8) and (-3.4,  1.0) .. (-2.4,  0.0)
  .. controls (-1.6, -0.8) and (-0.6, -0.6) .. ( 0,   -0.3)
  -- cycle;

\node[lbl, text=blue!60!black, font=\small\itshape] at (-1.5, 0.6) {$S$};
\node[lbl, font=\small\itshape] at (-3.8, 4.2) {$X$};

% --- Interior point x: ball fits ---
\coordinate (x1) at (-0.8, 2.2);
\draw[blue!70!black, dashed, line width=0.7pt] (x1) circle (0.85cm);
\draw[gray!60, line width=0.6pt, ->] (x1) -- ++(0.85, 0);
\node[sublbl, anchor=south] at ($(x1)+(0.42, 0.0)$) {$r$};
\node[dot, fill=blue!70!black] at (x1) {};
\node[lbl, anchor=south east] at ($(x1)+(-0.08, 0.10)$) {$x$};
\node[sublbl, anchor=south, align=center] at ($(x1)+(0, 0.92)$)
  {$B_r(x)\subseteq S$ \quad \textit{fits inside}};

% --- Near-boundary point z: ball leaks ---
\coordinate (z) at (1.35, 1.7);
\draw[red!60!black, dashed, line width=0.7pt] (z) circle (0.52cm);
\node[dot, fill=red!70!black] at (z) {};
\node[lbl, anchor=north west] at ($(z)+(0.08,-0.10)$) {$z$};
\coordinate (leak) at ($(z)+(0.46, 0.24)$);
\draw[red!70, line width=1.3pt]
  ($(leak)+(-0.09,-0.09)$) -- ($(leak)+(0.09, 0.09)$);
\draw[red!70, line width=1.3pt]
  ($(leak)+(-0.09, 0.09)$) -- ($(leak)+(0.09,-0.09)$);
\node[sublbl, anchor=west, align=left] at ($(z)+(0.65, 0.15)$)
  {$z\in S$, but every $B_r(z)$\\[1pt]meets $X\!\setminus\! S$};

% --- Footer ---
\node[lbl, align=center] at (-0.8, -1.15)
  {$S$ is \textbf{open}
   $\iff$
   $\forall x \in S,\;
    \exists r > 0,\;
    B_r(x) \subseteq S$};

\end{tikzpicture}
\captionof{figure}{The open set condition in a metric space $(X,d)$.
Interior point $x$: a ball $B_r(x)$ fits entirely inside $S$, witnessing
openness at $x$.
Near-boundary point $z$: every ball around $z$ leaks into $X\setminus S$,
so if such a point existed in $S$ then $S$ would fail to be open.
The dashed boundary of $S$ itself signals that $S$ is open (boundary excluded).}
\label{fig:open-set}
\end{center}

% ---------------------------------------------------------
\subsection{First Example: A Neither-Nor Set}
\label{sec:neither-nor-example}
% ---------------------------------------------------------

\begin{example}
Consider $A = [0,1) \subset \mathbb{R}$ with the standard metric.
Every ball around $0$ contains negative numbers, so $A$ is not open.
Every ball around $1$ contains points of $A$ just below $1$, so
$X \setminus A$ is not open either, hence $A$ is not closed.
The boundary is $\partial A = \{0, 1\}$; since $0 \in A$ but $1 \notin A$,
the set contains part of its boundary but not all.
Therefore $A$ is neither open nor closed.
\end{example}

\begin{remark*}[Openness and closedness are not about geometry]
Openness and closedness describe how a set interacts with
neighborhoods, not whether it has a visible boundary.
A set may have a boundary and be open, have a boundary and be closed,
or have a boundary and be neither.
What matters is the set's relationship to \emph{every} ball around
\emph{every} point.
\end{remark*}

% ---------------------------------------------------------
\subsection{First Propositions: Trivial Sets and Balls}
\label{sec:trivial-and-balls}
% ---------------------------------------------------------

\begin{proposition}[Trivial Open and Closed Sets]
\label{prop:trivial-open-closed}
Let $(X,d)$ be a metric space.
The empty set $\varnothing$ and the whole space $X$ are both open and closed.
\end{proposition}

\begin{proof}
$\varnothing$ is open because the defining condition
$\forall x \in U,\;\exists r>0,\;B_r(x)\subseteq U$
is vacuously true when $U = \varnothing$.

$X$ is open because for any $x \in X$ and any $r > 0$, the ball
$B_r(x) \subseteq X$ trivially.

Since $X \setminus \varnothing = X$ and $X \setminus X = \varnothing$ are both
open, both sets are also closed by definition.
\end{proof}

\begin{remark*}[Fully quantified form]
\begin{align*}
\varnothing\ \text{open}
  &\iff \forall x \in \varnothing,\;
         \exists r \in \mathbb{R}^{+},\;
         B_r(x) \subseteq \varnothing
  \quad\text{(vacuously true).} \\[4pt]
X\ \text{open}
  &\iff \forall x \in X,\;
         \exists r \in \mathbb{R}^{+},\;
         B_r(x) \subseteq X
  \quad\text{(true for any $r$).} \\[4pt]
\varnothing\ \text{closed}
  &\iff X \setminus \varnothing = X\ \text{is open.} \\[4pt]
X\ \text{closed}
  &\iff X \setminus X = \varnothing\ \text{is open.}
\end{align*}
\end{remark*}

\begin{remark*}[Consequence]
The topology axioms for open sets follow from
\hyperref[def:open-closed-set]{Definition~\ref*{def:open-closed-set}}:
arbitrary unions of open sets are open; finite intersections of open sets
are open.
By De Morgan, arbitrary intersections of closed sets are closed and finite
unions of closed sets are closed.
\end{remark*}

\begin{proposition}[\texorpdfstring{\hyperref[prf:balls-open-closed]{Balls are Open and Closed}}{Balls are Open and Closed}]
\label{prop:balls-open-closed}
Let $(X,d)$ be a metric space, $x \in X$, and $r > 0$.
\begin{enumerate}[label=(\roman*)]
  \item $B_r(x)$ is an open set.
  \item $\overline{B}_r(x)$ is a closed set.
\end{enumerate}
\end{proposition}

\begin{remark*}[Fully quantified form]
\begin{align*}
&\forall (X,d)\ \text{metric space},\;
 \forall x \in X,\;
 \forall r \in \mathbb{R}^{+}: \\
&\quad \text{(i)}\quad
  \forall y \in B_r(x),\;
  \exists \alpha \in \mathbb{R}^{+},\;
  B_\alpha(y) \subseteq B_r(x). \\
&\quad \text{(ii)}\quad
  X \setminus \overline{B}_r(x)\ \text{is open, i.e.}\;
  \forall y \notin \overline{B}_r(x),\;
  \exists \alpha \in \mathbb{R}^{+},\;
  B_\alpha(y) \cap \overline{B}_r(x) = \varnothing.
\end{align*}
\end{remark*}

\begin{remark*}[Intuition]
Both parts use the radial gap from
\hyperref[rem:interior-ball-identity]{\texttt{notes-balls.tex}}:

\smallskip
\begin{center}
\begin{tabular}{l l}
\toprule
\textbf{Case} & \textbf{Gap and safe radius} \\
\midrule
$y \in B_r(x)$ (inside open ball)
  & $\alpha := r - d(x,y) > 0$;\quad $B_\alpha(y) \subseteq B_r(x)$ \\[4pt]
$y \notin \overline{B}_r(x)$ (outside closed ball)
  & $\alpha := d(x,y) - r > 0$;\quad $B_\alpha(y) \cap \overline{B}_r(x) = \varnothing$ \\
\bottomrule
\end{tabular}
\end{center}

\smallskip\noindent
In both cases the triangle inequality converts the gap into a ball
containment.
\end{remark*}

% =========================================================
% TikZ diagram: Slack ball proof sketch for B_r(x) open
% Place after Proposition (Balls are Open and Closed).
% Requires: \usepackage{tikz}  \usetikzlibrary{calc}
%           \usepackage{caption}   (for \captionof)
% =========================================================

\begin{center}
\begin{tikzpicture}[
  dot/.style     = {circle, fill, inner sep=1.8pt},
  lbl/.style     = {font=\small},
  sublbl/.style  = {font=\scriptsize, text=gray!55!black}
]

\def\R{3.0}
\def\Dx{1.5}
\def\Dy{0.9}
\pgfmathsetmacro{\dxy}{sqrt(\Dx*\Dx+\Dy*\Dy)}
\pgfmathsetmacro{\al}{\R - \dxy}

% --- Outer ball B_r(x) ---
\fill[blue!8] (0,0) circle (\R cm);
\draw[blue!55!black, dashed, line width=1.0pt] (0,0) circle (\R cm);
\node[lbl, anchor=south west, text=blue!55!black]
  at ({-\R*0.707}, {\R*0.707}) {$B_r(x)$};

% --- Centre x ---
\node[dot, fill=black] at (0,0) {};
\node[lbl, anchor=north east] at (-0.08,-0.12) {$x$};
\draw[gray!60, ->, line width=0.6pt] (0,0) -- (0, \R);
\node[lbl, anchor=west] at (0.06, \R/2) {$r$};

% --- Point y ---
\coordinate (y) at (\Dx, \Dy);
\node[dot, fill=blue!70!black] at (y) {};
\node[lbl, anchor=south west] at ($(y)+(0.08, 0.08)$) {$y$};

% --- Slack ball B_alpha(y) ---
\fill[blue!20, opacity=0.7] (y) circle (\al cm);
\draw[blue!75!black, solid, line width=1.0pt] (y) circle (\al cm);
\node[sublbl, anchor=north, align=center] at ($(y)+(0,-\al-0.15)$)
  {$B_{\alpha}(y)$};

% --- d(x,y) line ---
\draw[gray!50, line width=0.5pt] (0,0) -- (y);
\node[sublbl, anchor=north, rotate=30] at (0.72, 0.28) {$d(x,y)$};

% --- alpha arrow ---
\pgfmathsetmacro{\nx}{\Dx/\dxy}
\pgfmathsetmacro{\ny}{\Dy/\dxy}
\pgfmathsetmacro{\ex}{\Dx + \al*\nx}
\pgfmathsetmacro{\ey}{\Dy + \al*\ny}
\draw[gray!60, ->, line width=0.6pt] (y) -- (\ex, \ey);
\node[sublbl, anchor=south west] at ($(\ex,\ey)+(0.06,0.06)$)
  {$\alpha = r - d(x,y)$};
\node[dot, fill=gray!50, inner sep=1.2pt] at (\ex, \ey) {};

% --- Annotation ---
\node[lbl, align=center] at ($(y)+(-1.8, -2.2)$)
  {$B_\alpha(y) \subseteq B_r(x)$\\[2pt]
   \textit{by the triangle inequality}};

% --- Key equation ---
\node[lbl, align=center] at (0, -3.6)
  {Choose $\alpha := r - d(x,y) > 0$.
   \quad
   If $d(y,z) < \alpha$,
   then
   $d(x,z) \le d(x,y) + d(y,z) < d(x,y) + \alpha = r$.};

\end{tikzpicture}
\captionof{figure}{Proof sketch that $B_r(x)$ is open
(Proposition~\ref{prop:balls-open-closed}(i)).
Given $y \in B_r(x)$, the slack $\alpha := r - d(x,y) > 0$ is the distance
from $y$ to the boundary of $B_r(x)$ along the ray through $x$ and $y$.
The triangle inequality then shows $B_\alpha(y) \subseteq B_r(x)$,
witnessing that $y$ is an interior point.}
\label{fig:slack-ball}
\end{center}

% ---------------------------------------------------------
\subsection{Interior, Exterior, and Boundary}
\label{sec:interior-exterior-boundary}
% ---------------------------------------------------------

\begin{tcolorbox}[colback=propbox, colframe=propborder, arc=2pt,
  left=6pt, right=6pt, top=4pt, bottom=4pt,
  title={\small\textbf{Definition (Interior)}},
  fonttitle=\small\bfseries]
\begin{definition}[Interior]
\label{def:interior}
Let $(X,d)$ be a metric space and let $A \subseteq X$.
The \emph{interior} of $A$ is
\[
  \operatorname{int}(A)
  \;:=\;
  \bigl\{x \in A : \exists r > 0 \text{ such that } B_r(x) \subseteq A\bigr\}.
\]
\end{definition}
\end{tcolorbox}

\begin{remark*}[Fully quantified form]
\[
\forall (X,d)\ \text{metric space},\;
\forall A \subseteq X,\;
\forall x \in X:\quad
x \in \operatorname{int}(A)
\;\iff\;
x \in A
\;\land\;
\exists r \in \mathbb{R}^{+},\;
\forall y \in X,\;
d(x,y) < r \Rightarrow y \in A.
\]
\end{remark*}

\begin{tcolorbox}[colback=propbox, colframe=propborder, arc=2pt,
  left=6pt, right=6pt, top=4pt, bottom=4pt,
  title={\small\textbf{Definition (Exterior)}},
  fonttitle=\small\bfseries]
\begin{definition}[Exterior]
\label{def:exterior}
Let $(X,d)$ be a metric space and let $A \subseteq X$.
The \emph{exterior} of $A$ is
\[
  \operatorname{ext}(A)
  \;:=\;
  \operatorname{int}(X \setminus A).
\]
Equivalently,
\[
  x \in \operatorname{ext}(A)
  \;\iff\;
  \exists r > 0 \text{ such that } B_r(x) \cap A = \varnothing.
\]
\end{definition}
\end{tcolorbox}

\begin{remark*}[Fully quantified form]
\[
\forall (X,d)\ \text{metric space},\;
\forall A \subseteq X,\;
\forall x \in X:\quad
x \in \operatorname{ext}(A)
\;\iff\;
\exists r \in \mathbb{R}^{+},\;
\forall y \in X,\;
d(x,y) < r \Rightarrow y \notin A.
\]
\end{remark*}

\begin{tcolorbox}[colback=propbox, colframe=propborder, arc=2pt,
  left=6pt, right=6pt, top=4pt, bottom=4pt,
  title={\small\textbf{Definition (Boundary)}},
  fonttitle=\small\bfseries]
\begin{definition}[Boundary]
\label{def:boundary}
Let $(X,d)$ be a metric space and let $A \subseteq X$.
The \emph{boundary} of $A$ is
\[
  \partial A
  \;:=\;
  \Bigl\{
    x \in X :
    \forall r > 0,\;
    B_r(x) \cap A \neq \varnothing
    \;\text{ and }\;
    B_r(x) \cap (X \setminus A) \neq \varnothing
  \Bigr\}.
\]
\end{definition}
\end{tcolorbox}

\begin{remark*}[Fully quantified form]
\[
\forall (X,d)\ \text{metric space},\;
\forall A \subseteq X,\;
\forall x \in X:\quad
x \in \partial A
\;\iff\;
\forall r \in \mathbb{R}^{+},\;
\bigl(\exists y \in A,\; d(x,y) < r\bigr)
\;\land\;
\bigl(\exists z \in X \setminus A,\; d(x,z) < r\bigr).
\]
\end{remark*}

\begin{remark*}[Intuition]
The three sets classify every point of $X$ by how it sits relative to $A$:

\smallskip
\begin{center}
\begin{tabular}{l l}
\toprule
\textbf{Region} & \textbf{Every sufficiently small ball \ldots} \\
\midrule
$\operatorname{int}(A)$ & lies entirely inside $A$ \\
$\operatorname{ext}(A)$ & lies entirely outside $A$ \\
$\partial A$            & straddles both $A$ and $X \setminus A$ \\
\bottomrule
\end{tabular}
\end{center}
\end{remark*}

\begin{remark*}[Diagram]
See the TikZ diagram in \texttt{diagrams.tex} for a visual illustration
of an interior point, boundary point, and exterior point with their
respective balls.
\end{remark*}

% ---------------------------------------------------------
\subsection{The Decomposition Theorem}
\label{sec:ibe-decomposition}
% ---------------------------------------------------------

\begin{theorem}[\texorpdfstring{\hyperref[prf:ibe-decomp]{Interior--Boundary--Exterior Decomposition}}{Interior-Boundary-Exterior Decomposition}]
\label{thm:ibe-decomp}
For any subset $A \subseteq X$,
\[
  X
  \;=\;
  \operatorname{int}(A)
  \;\dot{\cup}\;
  \partial A
  \;\dot{\cup}\;
  \operatorname{ext}(A),
\]
a disjoint union.
\end{theorem}

\begin{remark*}[Fully quantified form]
\begin{align*}
\forall (X,d)\ \text{metric space},\;
\forall A \subseteq X,\;
\forall x \in X:\quad
&\bigl(x \in \operatorname{int}(A)\bigr)
\;\lor\;
\bigl(x \in \partial A\bigr)
\;\lor\;
\bigl(x \in \operatorname{ext}(A)\bigr), \\[4pt]
\text{and these three cases are mutually exclusive:}\quad
&\operatorname{int}(A) \cap \partial A = \varnothing, \\
&\operatorname{int}(A) \cap \operatorname{ext}(A) = \varnothing, \\
&\partial A \cap \operatorname{ext}(A) = \varnothing.
\end{align*}
\end{remark*}

\begin{remark*}[Intuition]
Every point of the ambient space stands in exactly one of three
relations to $A$: firmly inside, firmly outside, or on the fence.
No point can belong to two regions: a ball around an interior point is
wholly inside $A$, which rules out both boundary and exterior membership.
\end{remark*}

\begin{remark*}[Consequence]
This decomposition immediately yields necessary and sufficient ball
conditions for openness and closedness, and for boundary ownership
(see below).
\end{remark*}

% ---------------------------------------------------------
\subsection{Boundary and Closure}
\label{sec:boundary-closure}
% ---------------------------------------------------------

\begin{tcolorbox}[colback=propbox, colframe=propborder, arc=2pt,
  left=6pt, right=6pt, top=4pt, bottom=4pt,
  title={\small\textbf{Definition (Adherent Point, Closure)}},
  fonttitle=\small\bfseries]
\begin{definition}[Adherent Point]
\label{def:adherent-point}
Let $(X,d)$ be a metric space and let $A \subseteq X$.
A point $x \in X$ is an \emph{adherent point} (or \emph{closure point}) of $A$ if
\[
  \forall r > 0,\quad B_r(x) \cap A \neq \varnothing.
\]
\end{definition}

\begin{definition}[Closure]
\label{def:closure}
The \emph{closure} of $A$, denoted $\overline{A}$, is the set of all adherent
points of $A$:
\[
  \overline{A}
  \;:=\;
  \{x \in X : \forall r > 0,\; B_r(x) \cap A \neq \varnothing\}.
\]
\end{definition}
\end{tcolorbox}

\begin{remark*}[Relationship to the boundary]
An adherent point of $A$ is any point that every ball around it touches $A$
--- it may lie inside $A$, or it may be a boundary point pressed up against $A$
from the outside.
Interior points and boundary points of $A$ are both adherent points;
exterior points are not.
\end{remark*}

\begin{theorem}[\texorpdfstring{\hyperref[prf:boundary-closure]{Boundary via Closure}}{Boundary via Closure}]
\label{thm:boundary-closure}
For any $A \subseteq X$,
\[
  \partial A
  \;=\;
  \overline{A} \cap \overline{(X \setminus A)},
  \qquad\qquad
  \overline{A} = A \cup \partial A.
\]
\end{theorem}

\begin{remark*}[Fully quantified form]
\begin{align*}
\forall (X,d)\ \text{metric space},\;
\forall A \subseteq X,\;
\forall x \in X:\quad
x \in \partial A
&\;\iff\;
\Bigl(
  \forall r \in \mathbb{R}^{+},\;
  B_r(x) \cap A \neq \varnothing
\Bigr)
\;\land\;
\Bigl(
  \forall r \in \mathbb{R}^{+},\;
  B_r(x) \cap (X \setminus A) \neq \varnothing
\Bigr). \\[6pt]
x \in \overline{A}
&\;\iff\;
x \in A \;\lor\; x \in \partial A \\
&\;\iff\;
\forall r \in \mathbb{R}^{+},\;
B_r(x) \cap A \neq \varnothing.
\end{align*}
\end{remark*}

\begin{remark*}[Intuition]
The boundary is precisely where both $A$ and its complement accumulate.
The closure sweeps up all of $A$ together with its boundary.
\end{remark*}

% ---------------------------------------------------------
\subsection{Openness and Closedness via the Boundary}
\label{sec:open-closed-boundary}
% ---------------------------------------------------------

\begin{proposition}[\texorpdfstring{\hyperref[prf:closed-owns-boundary]{Closed Sets Contain Their Boundary}}{Closed Sets Contain Their Boundary}]
\label{prop:closed-owns-boundary}
Let $F \subseteq X$. Then
\[
  F \text{ is closed}
  \;\iff\;
  \partial F \subseteq F.
\]
\end{proposition}

\begin{remark*}[Fully quantified form]
\begin{align*}
\forall (X,d)\ \text{metric space},\;
\forall F \subseteq X:\quad
F\ \text{closed}
&\;\iff\;
\forall x \in X,\;
\Bigl(
  \forall r \in \mathbb{R}^{+},\;
  B_r(x) \cap F \neq \varnothing
  \;\land\;
  B_r(x) \cap (X \setminus F) \neq \varnothing
\Bigr)
\Rightarrow x \in F. \\[6pt]
\forall (X,d)\ \text{metric space},\;
\forall U \subseteq X:\quad
U\ \text{open}
&\;\iff\;
\forall x \in U,\;
\exists r \in \mathbb{R}^{+},\;
B_r(x) \cap (X \setminus U) = \varnothing.
\end{align*}
\end{remark*}

\begin{remark*}[Intuition]
Closed sets \emph{own their boundary}: no boundary point can escape.
The dual statement is that $U$ is open if and only if
$U \cap \partial U = \varnothing$ --- open sets contain none of their
boundary. This follows directly from the decomposition theorem and the
definition of interior.
\end{remark*}

\begin{remark*}[Common error]
The converse inference ``sets that miss some boundary points are open''
is \emph{not} valid. A set can contain some boundary points and exclude
others, making it neither open nor closed.
\end{remark*}

% ---------------------------------------------------------
\subsection{Conceptual Summary}
\label{sec:open-closed-summary}
% ---------------------------------------------------------

\begin{center}
\[
  \boxed{\text{Open sets: room around every point (no boundary)}}
  \qquad
  \boxed{\text{Closed sets: own their boundary}}
\]
\end{center}

\begin{remark*}[Bounded sets]
\label{rem:bounded-set}
A set $A \subseteq X$ is \emph{bounded} if
$\exists x_0 \in X,\; \exists R > 0,\; \forall a \in A,\; d(a, x_0) \le R$.
Boundedness is independent of openness and closedness; a set may be any
combination of bounded/unbounded and open/closed.
\end{remark*}

\begin{remark*}[Forward reference]
The equivalent sequential characterisations --- closed sets are stable
under limits of sequences; open sets are stable under small perturbations
--- are developed in
\hyperref[sec:sequential-characterizations]{\S\ref*{sec:sequential-characterizations}}.
\end{remark*}